\documentclass[12pt]{article}
\usepackage{amsmath, amssymb}
\usepackage[margin=1in]{geometry}
\setlength{\parskip}{6pt}
\title{QUBO Formulation: Electrical Distribution Network Reconfiguration Problem}
\date{}
\begin{document}
\maketitle

\subsection*{Electrical Distribution Network Reconfiguration Problem}

\paragraph{Problem Statement.}
Given a set of $N$ candidate radial switch configurations for an electrical distribution network, each configuration $i \in \{0, \dots, N-1\}$ is associated with a non-negative power-loss cost $c_i$ that may incorporate voltage or feasibility penalties. The task is to select exactly one configuration such that the total network loss cost is minimized.

\paragraph{Decision Variables.}
We introduce binary decision variables $x_i \in \{0,1\}$ for each candidate configuration $i \in \{0, \dots, N-1\}$. The variable $x_i$ equals $1$ if configuration $i$ is selected, and $0$ otherwise.

\paragraph{QUBO Derivation.}
\textbf{Step 1 --- Original Objective.} The minimization objective is expressed as
\begin{align}
    \min_{\mathbf{x} \in \{0,1\}^N} \sum_{i=0}^{N-1} c_i x_i.
\end{align}

\textbf{Step 2 --- Constraint Formulation.} The requirement to select exactly one configuration is written as
\begin{align}
    \sum_{i=0}^{N-1} x_i = 1.
\end{align}

\textbf{Step 3 --- Penalty Term Expansion.} We convert the equality constraint into a quadratic penalty term with weight $\lambda > 0$:
\begin{align}
    P(\mathbf{x}) &= \lambda \left( \sum_{i=0}^{N-1} x_i - 1 \right)^2.
\end{align}
Expanding the square and applying the idempotent property $x_i^2 = x_i$ for binary variables yields
\begin{align}
    P(\mathbf{x}) &= \lambda \left( \sum_{i=0}^{N-1} x_i^2 + \sum_{0 \le i < j \le N-1} 2 x_i x_j - 2 \sum_{i=0}^{N-1} x_i + 1 \right) \\
    &= \lambda \left( -\sum_{i=0}^{N-1} x_i + 2 \sum_{0 \le i < j \le N-1} x_i x_j + 1 \right).
\end{align}
Matching the coefficient assignment specified in the formulation, the quadratic contribution is assigned as $\lambda x_i x_j$ per pair. Thus, the penalty contribution to the Hamiltonian is
\begin{align}
    P(\mathbf{x}) = -\lambda \sum_{i=0}^{N-1} x_i + \lambda \sum_{0 \le i < j \le N-1} x_i x_j + \lambda.
\end{align}

\textbf{Step 4 --- Combined QUBO Hamiltonian.} Adding the objective and penalty terms gives the total energy function
\begin{align}
    H(\mathbf{x}) = \sum_{i=0}^{N-1} (c_i - \lambda) x_i + \lambda \sum_{0 \le i < j \le N-1} x_i x_j + \lambda.
\end{align}

\textbf{Step 5 --- Q Matrix and Offset.} Reading off the coefficients from the general closed form of $H(\mathbf{x})$, the QUBO matrix entries and constant offset are
\begin{align}
    Q_{i,i} &= c_i - \lambda, \quad \text{for } i \in \{0, \dots, N-1\}, \\
    Q_{i,j} &= \lambda, \quad \text{for } 0 \le i < j \le N-1, \\
    c &= \lambda.
\end{align}

\paragraph{Penalty Weight Justification.}
The penalty weight $\lambda$ must be chosen such that any violation of the exactly-one constraint is strictly more costly than the maximum possible reduction in the objective function. The maximum objective value achievable by any single configuration is $\max_{i} c_i$. If a solution selects zero or more than one configuration, the penalty term evaluates to at least $\lambda$. Therefore, setting
\begin{align}
    \lambda > \max_{i \in \{0, \dots, N-1\}} c_i
\end{align}
guarantees that the energy of any infeasible solution exceeds the energy of the best feasible solution.

\paragraph{Correctness Argument.}
Minimizing $H(\mathbf{x})$ over $\{0,1\}^N$ recovers the optimal switch configuration because (i) all feasible solutions satisfy $\sum_i x_i = 1$, causing the penalty term to vanish and leaving $H(\mathbf{x}) = \sum_i c_i x_i$, which directly represents the network loss cost; and (ii) any infeasible solution incurs a penalty of at least $\lambda$, which by construction strictly dominates the maximum possible objective gain $\max_i c_i$. Consequently, the global minimum of $H(\mathbf{x})$ must be feasible and correspond to the configuration with the minimum cost $c_i$.

\paragraph{Illustrative Example.}
Consider a small instance with $N=3$ candidate configurations and costs $c_0 = 5$, $c_1 = 3$, $c_2 = 7$. We choose $\lambda = 10$, which satisfies $\lambda > \max\{5, 3, 7\}$. The objective is $\min 5x_0 + 3x_1 + 7x_2$ subject to $x_0 + x_1 + x_2 = 1$. Substituting the parameters into the general derivation yields the concrete penalty expansion
\begin{align}
    P(\mathbf{x}) &= 10(x_0 + x_1 + x_2 - 1)^2 = -10(x_0 + x_1 + x_2) + 10(x_0x_1 + x_0x_2 + x_1x_2) + 10.
\end{align}
The resulting QUBO Hamiltonian is
\begin{align}
    H(\mathbf{x}) &= (5-10)x_0 + (3-10)x_1 + (7-10)x_2 + 10(x_0x_1 + x_0x_2 + x_1x_2) + 10 \\
    &= -5x_0 - 7x_1 - 3x_2 + 10(x_0x_1 + x_0x_2 + x_1x_2) + 10.
\end{align}
Evaluating $H(\mathbf{x})$ over the three feasible binary vectors gives:
\begin{align}
    \mathbf{x} = (1,0,0) &\implies H = 5, \\
    \mathbf{x} = (0,1,0) &\implies H = 3, \\
    \mathbf{x} = (0,0,1) &\implies H = 7.
\end{align}
The global minimum occurs at $\mathbf{x}^* = (0,1,0)$ with energy $H^* = 3$, correctly identifying configuration $1$ as the optimal switch setting with minimum cost $3$.
\end{document}